\documentclass[aspectratio=169,10pt]{beamer}

\usetheme{Madrid}
\usecolortheme{whale}
\setbeamertemplate{navigation symbols}{}
\setbeamertemplate{caption}[numbered]
\setbeamertemplate{itemize items}[circle]

\usepackage{amsmath,amssymb}
\usepackage{booktabs}
\usepackage{listings}
\usepackage{xcolor}

\lstset{
  basicstyle=\ttfamily\scriptsize,
  keywordstyle=\color{blue!60!black},
  commentstyle=\color{green!35!black},
  stringstyle=\color{red!50!black},
  columns=fullflexible,
  keepspaces=true,
  showstringspaces=false,
  frame=single,
  framerule=0.3pt,
  rulecolor=\color{black!35},
  breaklines=true
}

\title{saturnCFD v1}
\subtitle{MPI and GPU finite-volume CFD with validated methodology}
\author{Ananth Narayan}
\institute{Indian Institute of Technology Madras}
\date{31 October 2026}

\begin{document}

\begin{frame}
  \titlepage
\end{frame}

\begin{frame}{Outline}
  \tableofcontents
\end{frame}

\section{Motivation and scope}

\begin{frame}{What saturnCFD is}
  \begin{block}{Goal}
    A single finite-volume code for scalar transport and incompressible flow,
    portable across CPU and GPU, parallel through MPI, driven entirely by YAML,
    and verified before it is trusted.
  \end{block}

  \begin{itemize}
    \item Structured hexahedral mesh imported from Gmsh \texttt{.msh}.
    \item 3D from the start; a one-cell-thick mesh is treated as 2D.
    \item Solvers: Burgers, heat/convection--diffusion, Stokes, incompressible
          Navier--Stokes.
    \item Kokkos for CPU/GPU; MPI for distributed memory.
    \item Every numerical choice is documented, cited and tested.
  \end{itemize}

  \vspace{0.4em}
  \small
  v0 is preserved on \texttt{v0-archive}. v1 is a deliberate rewrite, not a
  patch on the old AMR container.
\end{frame}

\begin{frame}{Why this order?}
  \begin{enumerate}
    \item Burgers: nonlinear convection, shocks, limiters, diffusion.
    \item Heat: scalar transport, Péclet effects, boundary layers.
    \item Stokes: vector elliptic operator, pressure coupling, linear solver.
    \item Incompressible Navier--Stokes: all of the above together.
  \end{enumerate}

  \begin{block}{Idea}
    Each PDE tests one more ingredient. A failure is localised before the next
    solver is added.
  \end{block}

  \small
  Compressible flow, AMR and two-phase flow are later milestones.
\end{frame}

\section{What is implemented}

\begin{frame}{Implemented substrate}
  \begin{itemize}
    \item Gmsh 2.2 reader; structured-hex classification and
          \texttt{checkMesh} report.
    \item Default \texttt{.vtu} + \texttt{.json}; optional
          \texttt{-writeVTK} patches and \texttt{.pvd}.
    \item yaml-cpp configuration with a typed schema and generated templates.
    \item SI-only validation: non-SI input is rejected, not converted.
    \item Single-dash CLI with help topics and shell completion.
    \item MPI Cartesian partition and 26-neighbour halo exchange.
    \item 1/2/4-rank MPI tests, regression tests, Doxygen HTML/PDF.
  \end{itemize}

  \begin{block}{Status}
    The mesh/config/comm slice is working. Kokkos fields and the FVM operators
    are the next implementation step.
  \end{block}
\end{frame}

\begin{frame}{Architecture}
  \begin{center}
  \small
  \begin{tabular}{ll}
    \toprule
    Layer & Responsibility \\
    \midrule
    \texttt{core} & types, SI units, errors, timer \\
    \texttt{config} & yaml-cpp, schema, validation, templates \\
    \texttt{mesh} & Gmsh, \texttt{checkMesh}, structured view, VTK \\
    \texttt{comm} & partition, ghost ownership, halo exchange \\
    \texttt{fv} & fields, gradient, divergence, Laplacian, fluxes \\
    \texttt{linalg} & MPI PETSc systems \\
    \texttt{simulation} & solver registry, time loop, diagnostics \\
    \texttt{solvers} & Burgers, heat, Stokes, incompressible \\
    \bottomrule
  \end{tabular}
  \end{center}

  \small
  A solver owns its equations. It does not own the mesh, the halo, the YAML
  parser or the linear algebra.
\end{frame}

\section{Mesh and configuration}

\begin{frame}{Structured mesh and \texttt{checkMesh}}
  \begin{itemize}
    \item Gmsh 2.2 ASCII: points, hexahedra, boundary physical groups.
    \item Classifier walks the cell dual graph and assigns a contiguous
          $(i,j,k)$.
    \item Rejects tetrahedra, non-manifold faces, negative volumes and
          non-contiguous hex meshes.
    \item Reports counts, patches, volume/area ranges, non-orthogonality,
          skewness, aspect ratio, structured flag and dimension.
  \end{itemize}

  \begin{block}{2D rule}
    A one-cell-thick 3D mesh with \texttt{empty} patches is 2D. One internal
    representation, one solver path.
  \end{block}

  \small
  Gmsh \cite{geuzaineRemacle2009}; OpenFOAM's \texttt{checkMesh} convention is a
  methodology reference.
\end{frame}

\begin{frame}[fragile]{YAML and SI units}
  \begin{columns}[T]
    \begin{column}{0.52\textwidth}
\begin{lstlisting}
solver: common
units:
  system: SI
mesh:
  file: channel.msh
time:
  end: 1.0
  cfl: 0.4
\end{lstlisting}
    \end{column}
    \begin{column}{0.44\textwidth}
      \begin{itemize}
        \item yaml-cpp parses; the schema validates.
        \item Unknown keys are errors.
        \item \texttt{-dumpConfig} generates the template from the same
              registration the solver reads.
        \item \texttt{-checkConfig} runs all validation before allocation.
      \end{itemize}
    \end{column}
  \end{columns}

  \vspace{0.5em}
  \begin{block}{SI only in v1}
    \texttt{units.system: SI} or no block is accepted. \texttt{CGS},
    \texttt{length: mm}, etc.\ are hard errors. A future units layer converts at
    the configuration boundary, never in a kernel.
  \end{block}
\end{frame}

\section{Parallel model}

\begin{frame}{MPI decomposition}
  \begin{center}
    \small
    \begin{tabular}{ll}
      \toprule
      global grid & $n_x \times n_y \times n_z$ \\
      process grid & $p_x \times p_y \times p_z$ \\
      local block & $[i_0,i_1]\times[j_0,j_1]\times[k_0,k_1]$ \\
      \bottomrule
    \end{tabular}
  \end{center}

  \begin{itemize}
    \item \texttt{MPI\_Dims\_create} chooses the process grid;
          \texttt{MPI\_Cart\_create} handles periodicity.
    \item 2D uses $p_z=1$; 1D uses $p_y=p_z=1$.
    \item Remainders are distributed one cell at a time from the low index.
    \item No METIS: a structured mesh does not need a graph partition.
  \end{itemize}

  \begin{block}{Why Cartesian blocks?}
    Deterministic neighbours, compact messages and a direct
    cell $\leftrightarrow$ rank mapping.
  \end{block}
\end{frame}

\begin{frame}{No shared cells: owner and ghost}
  \begin{columns}[T]
    \begin{column}{0.55\textwidth}
      \begin{itemize}
        \item Every global cell has \textbf{one owner rank}.
        \item Only the owner writes the valid value.
        \item The neighbour keeps a \textbf{read-only ghost copy}.
        \item Two ghost layers support MUSCL/TVD.
        \item 26-neighbour exchange: faces, edges, corners.
      \end{itemize}
    \end{column}
    \begin{column}{0.42\textwidth}
      \begin{block}{Shared face, not shared cell}
        Both ranks compute the same face flux from the same valid/ghost pair.
        Each rank updates only its own cell. No double counting.
      \end{block}
    \end{column}
  \end{columns}

  \vspace{0.8em}
  \small
  Halo order: local periodic copies $\rightarrow$ packed MPI $\rightarrow$
  physical boundary conditions.
\end{frame}

\begin{frame}{CPU and GPU execution}
  \begin{itemize}
    \item Kokkos selects the backend at \textbf{build} time:
          Serial/OpenMP or CUDA/HIP.
    \item CPU and GPU are separate build directories:
          \texttt{build-cpu} and \texttt{build-gpu}.
    \item \texttt{mpirun} launches MPI ranks; the compiled backend decides
          whether a rank uses a CPU or a GPU.
    \item Performance model: \textbf{one rank per GPU}.
    \item One GPU: \texttt{mpirun -np 1}. Four GPUs:
          \texttt{mpirun -np 4}.
    \item Halo first implementation: device pack $\rightarrow$ host MPI
          $\rightarrow$ device unpack.
  \end{itemize}

  \begin{block}{Case files are hardware-independent}
    CPU/GPU is not in \texttt{case.yaml}. Physics stays portable.
  \end{block}
\end{frame}

\section{Finite-volume method}

\begin{frame}{Collocated cell-centred FVM}
  \begin{itemize}
    \item Velocity and pressure at cell centroids.
    \item Face fluxes reconstructed from the two adjacent cells.
    \item Rhie--Chow correction suppresses checkerboard pressure.
    \item SIMPLE for steady Stokes/incompressible flow.
    \item PISO for transient incompressible flow.
  \end{itemize}

  \begin{block}{Why not MAC?}
    Staggered MAC gives exact discrete continuity, but it splits the velocity
    components onto different meshes and fights the vector form, \texttt{.msh}
    import and the 2D/3D rule. Collocated FV matches OpenFOAM and the default
    Basilisk/Nektar++ approach \cite{weller1998,cantwell2015}.
  \end{block}
\end{frame}

\begin{frame}{Governing equations}
  \small
  \begin{align*}
    \text{Burgers:}\quad
    &\frac{\partial u}{\partial t}
      + \nabla\cdot(u\mathbf{a}) = \nu\nabla^2 u + s \\[0.3em]
    \text{Heat:}\quad
    &\frac{\partial T}{\partial t}
      + \nabla\cdot(\mathbf{u}T)
      = \nabla\cdot(\alpha\nabla T) + S \\[0.3em]
    \text{Stokes:}\quad
    &-\nabla p + \mu\nabla^2\mathbf{u} = \mathbf{f},
      \qquad \nabla\cdot\mathbf{u}=0 \\[0.3em]
    \text{Incompressible NS:}\quad
    &\frac{\partial \mathbf{u}}{\partial t}
      + \nabla\cdot(\mathbf{u}\otimes\mathbf{u})
      = -\frac{1}{\rho}\nabla p + \nu\nabla^2\mathbf{u}
      + \frac{\mathbf{f}}{\rho}, \\
    &\nabla\cdot\mathbf{u}=0
  \end{align*}

  \vspace{-0.5em}
  \scriptsize
  Projection/pressure correction from Chorin \cite{chorin1968}; SIMPLE/PISO
  from Patankar \cite{patankar1980} and Issa \cite{issa1986}; Rhie--Chow
  \cite{rhieChow1983}.
\end{frame}

\begin{frame}{Vector form and linear algebra}
  \begin{block}{Discrete momentum}
    \[
      \mathbf{M}\frac{d\mathbf{u}}{dt}
      + \mathbf{C}(\mathbf{u})\mathbf{u}
      = -\mathbf{G}p + \mathbf{D}\mathbf{u} + \mathbf{f}
    \]
  \end{block}

  \begin{itemize}
    \item One vector unknown per cell: $u,v,w$.
    \item One sparsity pattern; block size 3 or three right-hand sides.
    \item PETSc: \texttt{Mat}, \texttt{Vec}, \texttt{KSP}, \texttt{PC}.
    \item Pressure Poisson: CG + GAMG/HYPRE; Neumann null space handled
          explicitly.
    \item No home-grown Krylov solver.
  \end{itemize}
\end{frame}

\section{Verification and validation}

\begin{frame}{Manufactured solutions: why}
  \begin{block}{Question MMS answers}
    Does the assembled discrete operator converge to the PDE at the expected
    order?
  \end{block}

  \begin{itemize}
    \item Unit tests check one function; regression checks one accepted result.
    \item MMS checks the full operator, including faces, ghosts, BCs,
          geometry and partition boundaries.
    \item It catches sign/coefficient errors, missed faces and halo bugs that
          are invisible on one rank.
  \end{itemize}

  \small
  Salari \& Knupp \cite{salari2000}; Roache \cite{roache2002}; Oberkampf \& Roy
  \cite{oberkampfRoy2010}.
\end{frame}

\begin{frame}{Manufactured solutions: method}
  Choose a smooth $\phi$, substitute into the PDE, and use the residual as the
  source:
  \[
    s = \mathcal{L}(\phi).
  \]

  For the periodic Laplacian:
  \[
    \phi=\sin(2\pi x)\sin(2\pi y)\sin(2\pi z),
    \qquad
    s=-3(2\pi)^2\phi.
  \]

  Refine $h$ and compute
  \[
    p \approx \frac{\log(e_{\mathrm{coarse}}/e_{\mathrm{fine}})}{\log 2}.
  \]

  \begin{block}{Expected orders}
    Central Laplacian/interpolation: 2. Upwind: 1. TVD on a smooth non-extremum:
    about 2, with local reduction at extrema.
  \end{block}
\end{frame}

\begin{frame}{Validation roadmap}
  \begin{center}
    \small
    \begin{tabular}{lll}
      \toprule
      Solver & First validation & Reference \\
      \midrule
      Burgers & viscous traveling wave & exact Cole--Hopf solution \\
      Burgers & inviscid shock & characteristics/conservation \\
      Heat & 1D boundary layer & exact solution \\
      Heat & 2D/3D MMS & manufactured order \\
      Stokes & Poiseuille/Couette & exact parabolic/linear \\
      Stokes & cavity $Re\to0$ & converged reference \\
      Incompressible & cavity $Re=100,1000$ & Ghia et al.\ \cite{ghia1982} \\
      Incompressible & Taylor--Green & energy decay/order \\
      \bottomrule
    \end{tabular}
  \end{center}

  \begin{block}{Rule}
    MMS and analytic tests before a physical benchmark. Observations are
    recorded, including failures and order reduction.
  \end{block}
\end{frame}

\section{Commands and tests}

\begin{frame}[fragile]{Build and run}
  \begin{columns}[T]
    \begin{column}{0.52\textwidth}
\begin{lstlisting}
# CPU build
cmake --preset cpu
cmake --build --preset cpu -j
cmake --install build-cpu \
  --prefix "$HOME/.local"

# GPU build
cmake --preset gpu
cmake --build --preset gpu -j
cmake --install build-gpu \
  --prefix "$HOME/.local"
\end{lstlisting}
    \end{column}
    \begin{column}{0.44\textwidth}
\begin{lstlisting}
# one installed command
mpirun -np 4 saturnCFD case.yaml
mpirun -np 1 saturnCFD case.yaml

# tests
ctest --preset default
ctest --preset regression
ctest --preset mpi
\end{lstlisting}
    \end{column}
  \end{columns}

  \vspace{0.4em}
  \small
  The installed binary carries the backend it was compiled with. CPU and GPU
  are separate build directories; install the one you want to run. Today
  \texttt{cmake --preset default} configures \texttt{build-cpu}; the explicit
  \texttt{cpu}/\texttt{gpu} presets and the Kokkos backend arrive with the FVM
  layer.
\end{frame}

\begin{frame}[fragile]{Help, clean and completion}
\begin{lstlisting}
saturnCFD -help
saturnCFD -help build
saturnCFD -help gpu
saturnCFD -help mpi
saturnCFD -help test
saturnCFD -help clean

source <(saturnCFD -completion bash)

cmake --build build-cpu --target clean
rm -rf build-cpu build-gpu
\end{lstlisting}

  \begin{block}{No memorisation}
    Test labels: \texttt{ctest --list-presets}. CLI: \texttt{saturnCFD -help}.
    Build: \texttt{-help build}. Clean: \texttt{-help clean}.
  \end{block}
\end{frame}

\begin{frame}{Current test status}
  \begin{center}
    \small
    \begin{tabular}{ll}
      \toprule
      Suite & Result \\
      \midrule
      CTest total & 16/16 passing \\
      MPI distribution \& halo & 1, 2, 4 ranks passing \\
      Mesh \texttt{checkMesh} & pass/reject fixtures passing \\
      Config SI-only & accept/reject fixtures passing \\
      CLI regression & help, outputs, completion, templates \\
      Doxygen & HTML and PDF build \\
      \bottomrule
    \end{tabular}
  \end{center}

  \begin{block}{What is not yet tested}
    Kokkos fields, FVM gradients/Laplacian, MMS orders, CPU/GPU parity and
    multi-GPU scaling. These are the next gates.
  \end{block}
\end{frame}

\section{Status and roadmap}

\begin{frame}{Status and next steps}
  \begin{columns}[T]
    \begin{column}{0.48\textwidth}
      \begin{block}{Done}
        \begin{itemize}
          \item Mesh + \texttt{checkMesh}
          \item YAML schema + SI validation
          \item CLI/help/completion/install
          \item MPI partition + 26-neighbour halo
          \item 1/2/4-rank tests
          \item Doxygen HTML/PDF
        \end{itemize}
      \end{block}
    \end{column}
    \begin{column}{0.48\textwidth}
      \begin{block}{Next}
        \begin{itemize}
          \item Kokkos \texttt{CellField} + device geometry
          \item Halo overload for fields
          \item gradient/divergence/Laplacian
          \item MMS periodic/Dirichlet/Neumann
          \item Burgers physics
          \item PETSc pressure solve
        \end{itemize}
      \end{block}
    \end{column}
  \end{columns}

  \vspace{0.6em}
  \begin{block}{Longer term}
    Stokes and incompressible NS; multi-GPU cluster scaling; AMR then two-phase
    flow; compressible flow.
  \end{block}
\end{frame}

\begin{frame}{Take-away}
  \begin{block}{What makes this code trustworthy}
    \begin{enumerate}
      \item A structured, validated mesh and a strict YAML/SI boundary.
      \item An ownership model with no shared cells and a tested halo.
      \item A single source for CPU/GPU through Kokkos, not two codes.
      \item Verification by manufactured solutions before benchmarks.
      \item A CLI and test layout that a user can discover, not memorise.
    \end{enumerate}
  \end{block}

  \vspace{0.8em}
  \centering
  \Large
  Questions and discussion
\end{frame}

\begin{frame}[allowframebreaks]{References}
  \bibliographystyle{plain}
  \bibliography{../references}
\end{frame}

\end{document}
